% ------------------------------------------------------------
% Chapter 11: Strong--Field Tests
% ------------------------------------------------------------

\chapter{Strong--Field Tests}
\label{chap:strong-field}

\section{Overview}

The observational success of general relativity in strong--field regimes
provides an essential benchmark for any alternative or emergent theory of
gravity. In this chapter we analyze \rsvpabbr{} predictions for compact
objects, gravitational waves, and the causal structure of strongly
lamphrodynamic regions. We pay particular attention to the question of
whether classical black holes exist in \rsvpabbr{}, and if so, how they
differ from their general--relativistic counterparts.

\section{Lamphrodynamic Compact Objects}

In \rsvpabbr{}, the structure of a compact object is determined not only by
its mass--energy density but also by the dynamics of the lamphrodynamic
fields $(\PhiField,\vField,\SField)$. The effective stress--energy tensor
contributing to the gravitational field is
\begin{equation}
  T_{ab}^{\mathrm{eff}}
   = T_{ab}^{\mathrm{matter}}
     + T_{ab}^{\mathrm{lam}},
\end{equation}
where $T_{ab}^{\mathrm{lam}}$ arises from lamphrodynamic gradients and
entropy flux. In spherical symmetry, the generalized Tolman--Oppenheimer--Volkoff
equations involve additional terms from $\SField$ and $\rho_{\mathrm{lam}}$,
leading to modified compactness bounds.

\section{Black Holes in RSVP?}

A significant question is whether classical black holes, with event horizons
and singularities, exist within \rsvpabbr{}. Lamphrodynamic transport and
entropy smoothing provide mechanisms that may prevent curvature singularities
from forming. Preliminary analysis suggests that strongly lamphrodynamic
regions undergo an entropy--driven ``decompression,'' reducing curvature and
potentially avoiding singularity formation.

\begin{conjecture}[No lamphrodynamic singularities]
\label{conj:no-singularity}
Under physically reasonable assumptions on entropy production and diffusion,
spacetime singularities are avoided in generic dynamical collapse scenarios.
\end{conjecture}

This conjecture requires further analysis; see \cref{chap:nonlinear-analysis}
for related discussion in the PDE context.

\section{Neutron Stars}

The structure equations for neutron stars in \rsvpabbr{} contain lamphrodynamic
pressure contributions. Depending on $\alpha_2$ and $\beta_2$ (see
\cref{chap:dark-sector}), the maximum neutron star mass and radius may differ
from general relativity, leading to potentially observable deviations in
neutron star mass--radius relationships.

\section{Gravitational Waves}

Gravitational waves in \rsvpabbr{} arise from metric perturbations influenced
by lamphrodynamic stresses. The propagation speed and polarization structure
may differ slightly from general relativity, depending on the coupling
parameters. Observations of gravitational waves from compact binaries provide
constraints on these couplings.

\begin{remark}[Wave speed constraint]
Observations from GW170817 constrain deviations in the gravitational wave
speed to be extremely small, limiting certain lamphrodynamic parameter
ranges.
\end{remark}

\section{Event Horizon Telescope Constraints}

Observations from the Event Horizon Telescope (EHT) provide tests of strong
gravity near supermassive compact objects. Lamphrodynamic modifications could
alter shadow size and photon ring structure. Current EHT data are consistent
with general relativity but do not rule out small lamphrodynamic effects.

\section{Discussion}

Strong--field observations place important constraints on lamphrodynamic
gravity, but the full parameter space remains open. Improved modeling of
compact objects and gravitational wave emission in \rsvpabbr{} is needed to
assess consistency with observations.

\clearpage
